% Copyright (c) 2026 Eugene Joseph M. Ragasa.
% SPDX-License-Identifier: Apache-2.0
% Author: Eugene Joseph M. Ragasa
% Formal proof completed: 2026-10-31
\documentclass[11pt]{article}
\usepackage[T1]{fontenc}\usepackage{amsmath,amssymb}\usepackage[margin=1in]{geometry}\usepackage[hidelinks]{hyperref}
\title{Proof Dossier: Pointwise-Unitary Fourier Support}\author{Eugene Joseph M. Ragasa}\date{Formal proof completed: 2026-10-31}
\begin{document}\maketitle
\section{Identity and status}
ID \texttt{BG-FOURIER-GAUGE-THM-002}; status \texttt{proved}. Lean declaration:
\par\noindent{\tiny\path{ProjectKoios.BandGauge.Fourier.Gauge.}}\\
{\tiny\path{pointwiseUnitaryConjugation_coefficients_supportedOn}}.
\section{Exact statement and proof}
If the Fourier coefficients of $U$ vanish outside $S_U$ and those of $A$ vanish outside $S_A$,
then
\[
 \operatorname{supp}\widehat{U^\dagger A U}
 \subseteq (-S_U+S_A)+S_U.
\]
Conjugate transpose and mode reversal place the left gauge factor on $-S_U$. Applying finite
convolution support addition first to that factor and $A$, then to the resulting convolution and the
right gauge factor, gives the declared support. The triple-convolution identity transfers this
support statement to the pointwise-conjugated field.
\section{Boundary}
The conclusion is inclusion in a supplied finite character inventory, not support equality. It does
not authenticate either support, infer quantitative localization or continuum decay, identify a
physical hopping range, or certify runtime gauge recovery or scientific adequacy.
\section{Trusted base and verification}
Checked with Lean \texttt{4.35.0} and Mathlib revision
\texttt{56a005855ea17e0b0958f73ec801ae140b0a1679}. Kernel axioms are
\texttt{propext}, \texttt{Classical.choice}, and \texttt{Quot.sound}; no admissions.
AI-assisted and kernel-checked; no independent expert review recorded.
\end{document}
